% ROZDZIAŁ 12 — do włączenia w main.tex
\chapter{Miary produktowe. Twierdzenia Tonellego i Fubiniego}
\label{ch:produkty-miary}

\section{Po co konstruujemy miarę produktową?}

Niech $(X,\mathcal F,\mu)$ i $(Y,\mathcal G,\nu)$ będą
przestrzeniami z miarami. Chcemy przypisać mierzalnemu
zbiorowi $E\subseteq X\times Y$ wielkość, która dla
prostokąta mierzalnego $A\times B$ wynosi $\mu(A)\nu(B)$.
Potrzebujemy następnie twierdzenia, które pozwoli obliczyć
całkę po $X\times Y$ jako dwie kolejne całki.

Już na etapie formułowania pojawiają się trzy pytania:
\emph{(i)} jaka jest odpowiednia $\sigma$-algebra
na $X\times Y$; \emph{(ii)} czy przypisanie
$A\times B\mapsto\mu(A)\nu(B)$ wyznacza miarę;
\emph{(iii)} kiedy można zmieniać kolejność całkowania.
Będziemy rozróżniać trzy odpowiedzi: konstrukcję miary
produktowej, twierdzenie Tonellego dla funkcji nieujemnych
i twierdzenie Fubiniego dla funkcji całkowalnych.

\section{Przestrzeń produktowa i przekroje}

Iloczyn $\sigma$-algebr $\mathcal F\otimes\mathcal G$
został zdefiniowany w Definicji~\ref{def:product-sigma}.
Zbiory $A\times B$ dla $A\in\mathcal F$ i $B\in\mathcal G$
nazywamy \emph{prostokątami mierzalnymi}.
Dla $E\subseteq X\times Y$ \emph{przekrojem pionowym}
przez $x\in X$ nazywamy $E_x=\{y\in Y:(x,y)\in E\}$,
a \emph{przekrojem poziomym} przez $y\in Y$ zbiór
$E^y=\{x\in X:(x,y)\in E\}$.
Dla funkcji $f:X\times Y\to\overline{\mathbb R}$ piszemy
$f_x(y)=f(x,y)$ i $f^y(x)=f(x,y)$.

\begin{lemat}[Mierzalność wszystkich przekrojów]
\label{lem:sections-measurable-12}
Niech $E\in\mathcal F\otimes\mathcal G$.
Wtedy $E_x\in\mathcal G$ dla każdego $x\in X$
oraz $E^y\in\mathcal F$ dla każdego $y\in Y$.
Jeżeli $f:(X\times Y,\mathcal F\otimes\mathcal G)
\to(\overline{\mathbb R},\mathcal B(\overline{\mathbb R}))$
jest mierzalna, to wszystkie funkcje $f_x$ oraz $f^y$
są mierzalne.
\end{lemat}
\begin{proof}
Ustalmy $x\in X$. Niech $\mathcal C_x$ oznacza rodzinę
tych zbiorów $E\subseteq X\times Y$, których przekrój
$E_x$ należy do $\mathcal G$. Ponieważ
$(X\times Y)_x=Y$, $(E^c)_x=Y\setminus E_x$ oraz
$(\bigcup_nE_n)_x=\bigcup_n(E_n)_x$, rodzina
$\mathcal C_x$ jest $\sigma$-algebrą. Dla prostokąta
$A\times B$ przekrój ma postać $B$, gdy $x\in A$,
oraz $\varnothing$, gdy $x\notin A$. Zatem
$\mathcal C_x$ zawiera wszystkie prostokąty i
$\mathcal F\otimes\mathcal G\subseteq\mathcal C_x$.
Argument dla $E^y$ jest symetryczny. Wreszcie
$\{y:f_x(y)\in H\}=(f^{-1}(H))_x$ dla każdego
zbioru borelowskiego $H\subseteq\overline{\mathbb R}$;
poprzednia część dowodu zapewnia jego mierzalność.
Analogicznie postępujemy z $f^y$.
\end{proof}

\begin{uwaga}[Nie odwracamy twierdzenia o przekrojach]
Z mierzalności wszystkich przekrojów nie wynika
mierzalność zbioru w $\mathcal F\otimes\mathcal G$.
Kontrprzykład z podzbiorem przekątnej podamy w dalszej
części rozdziału. Mierzalność przekrojów jest
\emph{konsekwencją} mierzalności produktowej, a nie
jej ogólnym kryterium.
\end{uwaga}

\section{Lemat o podwójnym liczeniu miary zbioru}

W tym miejscu nie mamy jeszcze zdefiniowanej miary na
$X\times Y$. Najpierw udowodnimy, że wielkości
$\int_X\nu(E_x)\,\mathrm d\mu(x)$ oraz
$\int_Y\mu(E^y)\,\mathrm d\nu(y)$ mają sens
i są równe, gdy obie miary są $\sigma$-skończone.
Argument jest podstawą konstrukcji miary produktowej.

\begin{lemat}[Przekroje na prostokącie skończonej miary]
\label{lem:finite-rectangle-sections-12}
Niech $C\in\mathcal F$, $D\in\mathcal G$ oraz
$\mu(C)<\infty$, $\nu(D)<\infty$. Jeżeli
$E\in\mathcal F\otimes\mathcal G$ i $E\subseteq C\times D$,
to funkcje $x\mapsto\nu(E_x)$ oraz
$y\mapsto\mu(E^y)$ są mierzalne i
\begin{equation}\label{eq:set-fubini-finite-12}
\int_C\nu(E_x)\,\mathrm d\mu(x)
=\int_D\mu(E^y)\,\mathrm d\nu(y).
\end{equation}
\end{lemat}
\begin{proof}
Rozpatrujemy przestrzeń mierzalną
$R=C\times D$ z $\sigma$-algebrą
$\mathcal H=(\mathcal F\otimes\mathcal G)|_R$.
Jest ona generowana przez prostokąty
$A\times B$ z $A\in\mathcal F|_C$ i
$B\in\mathcal G|_D$. Niech $\mathcal M$ będzie
rodziną tych $E\in\mathcal H$, dla których
obie funkcje przekrojowe są mierzalne oraz zachodzi
\eqref{eq:set-fubini-finite-12}.

\emph{Krok 1: algebra prostokątów.}
Dla $E=A\times B\subseteq R$ mamy
$\nu(E_x)=\nu(B)\chi_A(x)$ i
$\mu(E^y)=\mu(A)\chi_B(y)$. Obie całki
są równe $\mu(A)\nu(B)$.
Ta sama własność obowiązuje dla skończonej
\emph{rozłącznej} sumy prostokątów, z addytywności
miar i liniowości całki. Każdą skończoną sumę
prostokątów można przedstawić jako skończoną
rozłączną sumę prostokątów: wystarczy rozbić $C$
i $D$ na skończone rozłączne podziały wyznaczone
przez występujące pierwsze i drugie czynniki.
Zatem $\mathcal M$ zawiera algebrę $\mathcal A$
skończonych sum prostokątów.

\emph{Krok 2: ciągi rosnące.}
Jeżeli $E_n\in\mathcal M$ i $E_n\uparrow E$,
to dla każdego $x$ zachodzi $(E_n)_x\uparrow E_x$.
Ciągłość miary $\nu$ z dołu daje
$\nu(E_x)=\lim_n\nu((E_n)_x)$; zatem funkcja
$x\mapsto\nu(E_x)$ jest mierzalna jako granica
funkcji mierzalnych. Podobnie dla drugiego przekroju.
Twierdzenie Beppo Leviego zastosowane po obu stronach
\eqref{eq:set-fubini-finite-12} daje tę samą
równość dla $E$.

\emph{Krok 3: ciągi malejące.}
Jeżeli $E_n\in\mathcal M$ i $E_n\downarrow E$,
to $\nu((E_n)_x)\downarrow\nu(E_x)$, ponieważ
$\nu((E_1)_x)\le\nu(D)<\infty$. Ponadto funkcje
$\nu((E_n)_x)$ są zdominowane przez
$\nu(D)\chi_C(x)$, która ma skończoną
całkę $\mu(C)\nu(D)$ (w przypadku wartości zero
stosujemy zwykłą konwencję dla całki funkcji zerowej).
Twierdzenie o zbieżności zmajoryzowanej daje zbieżność
całek; analogicznie z drugiej strony, z dominantą
$\mu(C)\chi_D$. Stąd równość dla $E$.

Z kroków 2--3 wynika, że $\mathcal M$ jest klasą
monotoniczną. Zawiera algebrę $\mathcal A$, więc
na mocy twierdzenia o klasie monotonicznej
$\mathcal H=\sigma(\mathcal A)\subseteq\mathcal M$.
To kończy dowód bez użycia Fubiniego lub Tonellego.
\end{proof}

\begin{twierdzenie}[Tożsamość przekrojowa przy $\sigma$-skończoności]
\label{thm:section-identity-12}
Niech $\mu$ i $\nu$ będą $\sigma$-skończone.
Dla każdego $E\in\mathcal F\otimes\mathcal G$
funkcje $x\mapsto\nu(E_x)$ i $y\mapsto\mu(E^y)$
są mierzalne oraz
\begin{equation}\label{eq:section-identity-12}
\int_X\nu(E_x)\,\mathrm d\mu(x)
=\int_Y\mu(E^y)\,\mathrm d\nu(y)
\quad\text{w }[0,\infty].
\end{equation}
\end{twierdzenie}
\begin{proof}
Wybieramy rosnące wyczerpania
$X_n\uparrow X$ i $Y_n\uparrow Y$ zbiorami mierzalnymi
o skończonych miarach. Można je otrzymać, zastępując
zbiory z dowolnych wyczerpań ich skończonymi sumami.
Niech $R_n=X_n\times Y_n$ oraz $E_n=E\cap R_n$.
Lemat~\ref{lem:finite-rectangle-sections-12}
zastosowany do $E_n$ zapewnia mierzalność obu
funkcji przekrojowych dla $E_n$ i równość odpowiednich
całek. Ponieważ $E_n\uparrow E$, dla każdego $x$
zachodzi $(E_n)_x\uparrow E_x$, a dla każdego $y$
zachodzi $(E_n)^y\uparrow E^y$. Z ciągłości miary
z dołu wynika
$\nu(E_x)=\lim_n\nu((E_n)_x)$ i
$\mu(E^y)=\lim_n\mu((E_n)^y)$. Granice są
mierzalne. Stosując dwukrotnie twierdzenie
Beppo Leviego do całek po $X$ i $Y$, uzyskujemy
\eqref{eq:section-identity-12}.
\end{proof}

\section{Istnienie i jednoznaczność miary produktowej}

\begin{twierdzenie}[Konstrukcja miary produktowej]
\label{thm:product-measure-12}
Niech $(X,\mathcal F,\mu)$ i $(Y,\mathcal G,\nu)$
będą przestrzeniami z miarami $\sigma$-skończonymi.
Istnieje dokładnie jedna miara $\pi$ na
$(X\times Y,\mathcal F\otimes\mathcal G)$ taka, że
\begin{equation}\label{eq:rect-product-12}
\pi(A\times B)=\mu(A)\nu(B),
\qquad A\in\mathcal F,\quad B\in\mathcal G,
\end{equation}
przy umowie $0\cdot\infty=\infty\cdot0=0$.
Oznaczamy ją przez $\mu\otimes\nu$.
Dla każdego $E\in\mathcal F\otimes\mathcal G$ zachodzi
\begin{equation}\label{eq:measure-section-12}
(\mu\otimes\nu)(E)=
\int_X\nu(E_x)\,\mathrm d\mu(x)
=\int_Y\mu(E^y)\,\mathrm d\nu(y).
\end{equation}
\end{twierdzenie}
\begin{proof}
\emph{Istnienie.}
Definiujemy $\pi(E)$ jako pierwszą całkę
w \eqref{eq:measure-section-12}; na mocy
\ref{thm:section-identity-12} całkowana funkcja jest
mierzalna i można też użyć drugiej całki.
Oczywiście $\pi(\varnothing)=0$.
Jeśli $(E_j)_{j\ge1}$ jest ciągiem parami
rozłącznych zbiorów w $\mathcal F\otimes\mathcal G$,
to dla każdego $x$ przekroje $(E_j)_x$ są parami
rozłączne i
\[
\nu\left(\left(\bigcup_{j\ge1}E_j\right)_x\right)
=\sum_{j\ge1}\nu((E_j)_x).
\]
Stosując twierdzenie Beppo Leviego do sum częściowych,
ostatecznie otrzymujemy
$\pi(\bigcup_jE_j)=\sum_j\pi(E_j)$.
Zatem $\pi$ jest miarą.
Dla $E=A\times B$ funkcja przekrojowa jest równa
$\nu(B)\chi_A$ (w rozumieniu funkcji przyjmującej
wartość zero poza $A$), więc
$\pi(A\times B)=\mu(A)\nu(B)$ także w przypadkach
$\mu(A)=0$ lub $\nu(B)=0$.

\emph{Jednoznaczność.}
Niech $\rho$ będzie inną miarą spełniającą
\eqref{eq:rect-product-12}. Z wyczerpań użytych
w poprzednim twierdzeniu wybieramy prostokąty
$R_n=X_n\times Y_n\uparrow X\times Y$, gdzie
$\mu(X_n)\nu(Y_n)<\infty$.
Na $\sigma$-algebrze śladu na $R_n$ miary
$\rho$ i $\pi$ są skończone i zgadzają się na
prostokątach $A\times B\subseteq R_n$.
Przecięcia dwóch prostokątów są prostokątami,
a prostokąty generują $\sigma$-algebrę śladu;
twierdzenie $\pi$--$\lambda$ z rozdziału 1 daje
równość miar na całym śladzie.
W szczególności dla dowolnego
$E\in\mathcal F\otimes\mathcal G$ mamy
$\rho(E\cap R_n)=\pi(E\cap R_n)$.
Przechodząc do granicy $n\to\infty$ przez
ciągłość miar z dołu, dostajemy $\rho(E)=\pi(E)$.
\end{proof}

\begin{uwaga}[Co daje $\sigma$-skończoność?]
W powyższym dowodzie jest potrzebna do
wyczerpania przestrzeni prostokątami skończonej
miary, a w konsekwencji do jednoznaczności.
Bez $\sigma$-skończoności równość iterowanych całek
i jednoznaczność rozszerzenia nie muszą zachodzić.
Nie należy traktować jej jako dekoracyjnego założenia.
\end{uwaga}

\section{Twierdzenie Tonellego}

\begin{twierdzenie}[Tonelli]
\label{thm:tonelli-12}
Niech $\mu$ i $\nu$ będą miarami $\sigma$-skończonymi,
a $f:X\times Y\to[0,\infty]$ funkcją
$\mathcal F\otimes\mathcal G$-mierzalną.
Wtedy funkcje
\[
x\longmapsto\int_Y f(x,y)\,\mathrm d\nu(y),
\qquad
y\longmapsto\int_X f(x,y)\,\mathrm d\mu(x)
\]
są mierzalne, a trzy całki są równe
(wartość $+\infty$ jest dozwolona):
\begin{equation}\label{eq:tonelli-12}
\int_{X\times Y}f(x,y)\,\mathrm d(\mu\otimes\nu)(x,y)
=\int_X\left(\int_Yf(x,y)\,\mathrm d\nu(y)\right)
\mathrm d\mu(x)
=\int_Y\left(\int_Xf(x,y)\,\mathrm d\mu(x)\right)
\mathrm d\nu(y).
\end{equation}
\end{twierdzenie}
\begin{proof}
\emph{Krok 1: funkcje charakterystyczne.}
Dla $f=\chi_E$, $E\in\mathcal F\otimes\mathcal G$,
mamy $\int_Yf(x,y)\,\mathrm d\nu(y)=\nu(E_x)$.
Obie równości w \eqref{eq:tonelli-12} są dokładnie
wzorem \eqref{eq:measure-section-12}.

\emph{Krok 2: funkcje proste nieujemne.}
Jeśli $s=\sum_{j=1}^m a_j\chi_{E_j}$,
gdzie $a_j\ge0$, $E_j$ są mierzalne i parami
rozłączne, to liniowość całki skończonej sumy
funkcji nieujemnych daje
$\int_Y s(x,y)\,\mathrm d\nu(y)
=\sum_j a_j\nu((E_j)_x)$.
Wyrażenie po prawej jest mierzalne w $x$.
Całkując po $X$ i używając kroku 1,
otrzymujemy pierwszą równość dla $s$;
analogicznie postępujemy z odwrotną kolejnością.

\emph{Krok 3: dowolne funkcje nieujemne.}
Na mocy konstrukcji funkcji prostych z rozdziału 6
istnieje ciąg $s_n\uparrow f$ nieujemnych,
prostych i mierzalnych na $X\times Y$.
Dla każdego $x$ mamy $(s_n)_x\uparrow f_x$,
więc twierdzenie Beppo Leviego względem $\nu$
daje $\int_Y(s_n)_x\,\mathrm d\nu(y)\uparrow
\int_Yf_x\,\mathrm d\nu(y)$. Granica funkcji
mierzalnych jest mierzalna. Ponownie stosujemy
Beppo Leviego, tym razem względem $\mu$:
\[
\int_X\left(\int_Y f_x\,\mathrm d\nu(y)\right)\mathrm d\mu(x)
=\lim_n\int_X\left(\int_Y (s_n)_x\,\mathrm d\nu(y)\right)\mathrm d\mu(x)
=\lim_n\int_{X\times Y}s_n\,\mathrm d(\mu\otimes\nu)(x,y).
\]
Ostatnia granica jest równa
$\int_{X\times Y}f\,\mathrm d(\mu\otimes\nu)(x,y)$,
jeszcze raz z twierdzenia Beppo Leviego.
Argument z całkowaniem najpierw względem $\mu$
jest identyczny.
\end{proof}

\begin{wniosek}[Kryterium nieujemnej całkowalności]
\label{cor:tonelli-integrability-12}
W sytuacji twierdzenia~\ref{thm:tonelli-12}
niech $f\ge0$. Funkcja $f$ ma skończoną całkę
względem $\mu\otimes\nu$ wtedy i tylko wtedy,
gdy dowolna (a zatem każda) z iterowanych całek
w \eqref{eq:tonelli-12} jest skończona.
\end{wniosek}
\begin{proof}
Wszystkie trzy całki są równe jako elementy
$[0,\infty]$. Skończoność jednej jest zatem
równoważna skończoności pozostałych.
\end{proof}

\section{Twierdzenie Fubiniego dla funkcji całkowalnych}

\begin{twierdzenie}[Fubini]
\label{thm:fubini-12}
Niech $\mu$ i $\nu$ będą miarami $\sigma$-skończonymi,
a $f:X\times Y\to\mathbb R$ lub $\mathbb C$ funkcją
$\mathcal F\otimes\mathcal G$-mierzalną taką, że
\begin{equation}\label{eq:abs-fubini-12}
\int_{X\times Y}|f(x,y)|\,\mathrm d(\mu\otimes\nu)(x,y)<\infty.
\end{equation}
Wtedy:
\begin{enumerate}[label=\textup{(\roman*)}]
\item dla $\mu$-prawie każdego $x$ przekrój $f_x$ jest
      całkowalny względem $\nu$; analogicznie dla
      $\nu$-prawie każdego $y$ przekrój $f^y$ jest
      całkowalny względem $\mu$;
\item przyjmując wartość zero na odpowiednich
      mierzalnych zbiorach wyjątkowych otrzymujemy
      całkowalne funkcje
      $F(x)=\int_Y f(x,y)\,\mathrm d\nu(y)$ i
      $G(y)=\int_X f(x,y)\,\mathrm d\mu(x)$;
\item zachodzi
\begin{equation}\label{eq:fubini-12}
\int_{X\times Y}f(x,y)\,\mathrm d(\mu\otimes\nu)(x,y)
=\int_XF(x)\,\mathrm d\mu(x)
=\int_YG(y)\,\mathrm d\nu(y).
\end{equation}
\end{enumerate}
\end{twierdzenie}
\begin{proof}
\emph{Krok 1: całkowalność przekrojów prawie wszędzie.}
Stosując twierdzenie Tonellego do $|f|$,
otrzymujemy
\[
\int_X\left(\int_Y|f(x,y)|\,\mathrm d\nu(y)\right)
\mathrm d\mu(x)
=\int_{X\times Y}|f|\,\mathrm d(\mu\otimes\nu)(x,y)<\infty.
\]
Funkcja nieujemna $H(x)=\int_Y|f(x,y)|\,\mathrm d\nu(y)$
jest mierzalna i ma skończoną całkę. Zatem
$N=\{x:H(x)=\infty\}$ jest mierzalnym zbiorem
miary zero. Dla $x\notin N$ przekrój $f_x$
jest całkowalny; analogicznie otrzymujemy zbiór
wyjątkowy $M\subseteq Y$ miary zero.

\emph{Krok 2: mierzalność całek przekrojowych.}
Dla funkcji rzeczywistej zapisujemy $f=f^+-f^-$.
Funkcje $x\mapsto\int_Y f^+(x,y)\,\mathrm d\nu(y)$
i $x\mapsto\int_Y f^-(x,y)\,\mathrm d\nu(y)$
są mierzalne z Tonellego. Dla $x\notin N$
obie wartości są skończone, więc ich różnica
jest dobrze określona i mierzalna. Przedłużając
różnicę zerem na $N$, otrzymujemy mierzalną
funkcję $F$. Dla funkcji zespolonej powtarzamy
tę konstrukcję osobno dla części rzeczywistej
i urojonej. Analogicznie postępujemy z $G$.

\emph{Krok 3: całkowalność i równość całek.}
Dla $x\notin N$ nierówność trójkąta dla całki
daje $|F(x)|\le H(x)$, a na $N$ funkcja $F$
jest równa zero. Zatem $F\in L^1(\mu)$;
podobnie $G\in L^1(\nu)$.
Do funkcji $f^+$ i $f^-$ stosujemy Tonellego,
a następnie odejmujemy skończone liczby;
uzyskujemy pierwszą i drugą równość
\eqref{eq:fubini-12}. Dla funkcji zespolonej
stosujemy ten argument do obu części składowych.
\end{proof}

\begin{wniosek}[Kryterium praktyczne]
\label{cor:practical-fubini-12}
Jeśli dla mierzalnej funkcji $f:X\times Y\to\mathbb C$
możemy wykazać
$\int_X\bigl(\int_Y |f(x,y)|\,\mathrm d\nu(y)\bigr)
\mathrm d\mu(x)<\infty$, to $f$ jest całkowalna
na przestrzeni produktowej i wolno zamienić kolejność
całkowania w \eqref{eq:fubini-12}.
\end{wniosek}
\begin{proof}
Twierdzenie Tonellego zastosowane do $|f|$ zapewnia
równość podanej iterowanej całki i całki
$\int_{X\times Y}|f|\,\mathrm d(\mu\otimes\nu)(x,y)$.
Skończoność tej ostatniej oznacza spełnienie
warunków Fubiniego.
\end{proof}

\section{Przykłady rachunkowe}

\begin{przyklad}[Pole trójkąta bez całkowania po powierzchni]
Niech $T=\{(x,y)\in[0,1]^2:0\le y\le x\}$.
Przekrój pionowy $T_x$ jest przedziałem $[0,x]$
dla $x\in[0,1]$, a poza tym zbiorem jest pusty.
Dlatego Tonelli daje
\[
(\lambda_1\otimes\lambda_1)(T)=\int_{[0,1]}\left(
\int_{[0,1]}\chi_{\{y\le x\}}\,\mathrm d\lambda_1(y)
\right)\mathrm d\lambda_1(x)
=\int_{[0,1]}x\,\mathrm d\lambda_1(x)=\frac12.
\]
Z twierdzenia~\ref{thm:lebesgue-products-12} poniżej wynika także
$\lambda_2(T)=(\lambda_1\otimes\lambda_1)(T)=1/2$.
W odwrotnej kolejności przekrój poziomy ma długość
$1-y$, a całka wynosi
$\int_{[0,1]}(1-y)\,\mathrm d\lambda_1(y)=1/2$.
Nawet w tym elementarnym przykładzie widać dwie
różne parametryzacje tej samej miary.
\end{przyklad}

\begin{przyklad}[Całka z nieograniczoną funkcją nieujemną]
Na $(0,1)^2$ określmy $f(x,y)=\frac{1}{x+y}$.
Funkcja jest mierzalna i nieujemna; dlatego wolno
najpierw całkować po $y$, mimo że jeszcze nie
wiemy, czy podwójna całka jest skończona:
\[
\int_{(0,1)^2}\frac{1}{x+y}\,\mathrm d\lambda_2(x,y)
=\int_{(0,1)}\bigl(\log(1+x)-\log x\bigr)
\,\mathrm d\lambda_1(x)=2\log2.
\]
Ostatnia równość wynika z elementarnego
całkowania funkcji $\log(1+x)$ i $\log x$
na $(0,1)$ oraz przejścia do granicy przy zerze.
\end{przyklad}

\begin{przyklad}[Funkcja zmieniająca znak]
Niech $f(x,y)=e^{-x-y}(x-y)$ na $(0,\infty)^2$.
Ponieważ $|x-y|\le x+y$, Tonelli oraz wzory
$\int_{(0,\infty)}e^{-x}\,\mathrm d\lambda_1(x)=1$ i
$\int_{(0,\infty)}xe^{-x}\,\mathrm d\lambda_1(x)=1$
z rozdziału~\ref{ch:parametry-calki} dają
\begin{align*}
\int_{(0,\infty)^2}|f|\,
\mathrm d(\lambda_1\otimes\lambda_1)
&\le\int_{(0,\infty)^2}(x+y)e^{-x-y}
\,\mathrm d(\lambda_1\otimes\lambda_1)\\
&=1\cdot1+1\cdot1=2.
\end{align*}
Zatem wolno użyć Fubiniego. Zamiana $x$ z $y$
zmienia znak całkowanej funkcji i zachowuje
miarę produktową, a więc
$\int_{(0,\infty)^2}f\,\mathrm d(\lambda_1\otimes\lambda_1)=0$.
\end{przyklad}

\section{Dlaczego założenia twierdzeń są ważne?}

\begin{przyklad}[Dwie iterowane całki, dwa różne wyniki]
\label{ex:nonabs-fubini-12}
Niech $X=Y=\mathbb N$ z $\sigma$-algebrą potęgową
i miarą liczącą $\kappa$. Definiujemy
\[
a_{m,n}=\begin{cases}
1,&n=m,\\
-1,&n=m+1,\\
0,&\text{w pozostałych przypadkach}.
\end{cases}
\]
Dla każdego $m$ szereg $\sum_{n=1}^\infty a_{m,n}$
ma tylko dwa niezerowe wyrazy i wynosi zero. Natomiast
dla $n=1$ szereg $\sum_{m=1}^\infty a_{m,1}$
ma sumę $1$, a dla $n\ge2$ sumę zero.
Wobec tego
\[
\sum_{m=1}^\infty\left(\sum_{n=1}^\infty a_{m,n}\right)=0,
\qquad
\sum_{n=1}^\infty\left(\sum_{m=1}^\infty a_{m,n}\right)=1.
\]
Nie ma sprzeczności z Fubinim, ponieważ
$\sum_{m,n}|a_{m,n}|=\infty$.
W rzeczywistości $a^+$ i $a^-$ mają nieskończoną
całkę względem $\kappa\otimes\kappa$, więc
całka podpisana po całej przestrzeni nie istnieje.
\end{przyklad}

\begin{przyklad}[Bez $\sigma$-skończoności Tonelli może zawieść]
\label{ex:non-sigma-tonelli-12}
Niech $X=Y=[0,1]$, przy czym na $X$ bierzemy
$\mathcal F=\mathcal P([0,1])$ i miarę liczącą $\kappa$,
a na $Y$ miarę Lebesgue'a $\lambda_1$ ze zwykłą
$\sigma$-algebrą. Miara $\kappa$ nie jest
$\sigma$-skończona, gdyż jej zbiory skończonej
miary są skończone, a przeliczalna suma zbiorów
skończonych jest przeliczalna.
Zbiór przekątnej $D=\{(x,y):x=y\}$ jest
mierzalny w $\mathcal F\otimes\mathcal L_1$,
gdyż jest domknięty w zwykłej topologii produktu,
a $\mathcal B([0,1]^2)\subseteq\mathcal P([0,1])\otimes\mathcal L_1$:
każdy prostokąt otwarty należy do tej $\sigma$-algebry,
a otwarte podzbiory $[0,1]^2$ są przeliczalnymi sumami
takich prostokątów.
Dla każdego $x$ przekrój $D_x=\{x\}$ ma
$\lambda_1$-miarę zero, natomiast dla każdego $y$
przekrój $D^y=\{y\}$ ma miarę liczącą $1$.
Dlatego
\[
\int_X\left(\int_Y\chi_D(x,y)\,
\mathrm d\lambda_1(y)\right)\mathrm d\kappa(x)=0,
\qquad
\int_Y\left(\int_X\chi_D(x,y)\,
\mathrm d\kappa(x)\right)\mathrm d\lambda_1(y)=1.
\]
Nie można bez dodatkowych warunków konstruować
jednej miary produktowej przez dowolną z dwóch
kolejności całkowania. Ten przykład pokazuje
dokładnie miejsce, w którym używaliśmy
$\sigma$-skończoności w dowodzie.
\end{przyklad}

\section{Produkt miar Lebesgue'a i problem zupełności}

\begin{twierdzenie}[Produkt miar Lebesgue'a]
\label{thm:lebesgue-products-12}
Niech $m,n\ge1$. Miara $\lambda_m\otimes\lambda_n$
na $\mathcal L_m\otimes\mathcal L_n$ jest ograniczeniem
miary Lebesgue'a $\lambda_{m+n}$ do tej
$\sigma$-algebry. Ponadto uzupełnienie
$(\mathcal L_m\otimes\mathcal L_n,
\lambda_m\otimes\lambda_n)$ daje dokładnie
$(\mathcal L_{m+n},\lambda_{m+n})$.
W ogólności $\mathcal L_m\otimes\mathcal L_n$
nie jest zupełna.
\end{twierdzenie}
\begin{proof}
\emph{Krok 1: relacja między $\sigma$-algebrami.}
Dla $A\in\mathcal L_m$ istnieje zbiór borelowski
$B\subseteq\mathbb R^m$ taki, że
$A\mathbin\triangle B\subseteq N$ dla pewnego
borelowskiego zbioru zerowego $N$.
Analogicznie postępujemy z $C\in\mathcal L_n$.
Prostokąt $A\times C$ różni się od
borelowskiego prostokąta $B\times D$ co najwyżej
na podzbiorze
$(N\times\mathbb R^n)\cup(\mathbb R^m\times M)$,
gdzie $N$ i $M$ są borelowskie i mają miarę zero.
Zbiory te są $\lambda_{m+n}$-zerowe:
po przecięciu z ograniczonymi prostopadłościanami
wynika to z pokrywania małymi prostokątami
lub ze znanej już konstrukcji miary Lebesgue'a,
a następnie z przeliczalnej sumy po rosnących
prostopadłościanach. Stąd każdy prostokąt
$A\times C$ jest $\lambda_{m+n}$-mierzalny,
i $\mathcal L_m\otimes\mathcal L_n\subseteq\mathcal L_{m+n}$.

\emph{Krok 2: zgodność miar.}
Niech najpierw $Q\subseteq\mathbb R^m$ będzie ograniczonym
półotwartym prostopadłościanem. Dla borelowskich
$D\subseteq\mathbb R^n$ obie funkcje
$D\mapsto\lambda_{m+n}(Q\times D)$ oraz
$D\mapsto\lambda_m(Q)\lambda_n(D)$ są miarami
$\sigma$-skończonymi na $\mathcal B(\mathbb R^n)$;
zgadzają się na półotwartych prostopadłościanach
w $\mathbb R^n$, ponieważ ich produkty z $Q$ są
prostopadłościanami w $\mathbb R^{m+n}$.
Twierdzenie o jednoznaczności miar z rozdziału 2
daje równość dla każdego borelowskiego $D$.

Teraz ustalmy borelowski $D\subseteq\mathbb R^n$
o $\lambda_n(D)<\infty$. Obie funkcje
$B\mapsto\lambda_{m+n}(B\times D)$ oraz
$B\mapsto\lambda_n(D)\lambda_m(B)$ są
$\sigma$-skończonymi miarami borelowskimi na
$\mathbb R^m$ (gdy $\lambda_n(D)=0$, obie są
zerowe). Dzięki poprzedniemu akapitowi zgadzają
się na ograniczonych półotwartych prostopadłościanach
$B=Q$. Ponowne zastosowanie jednoznaczności
daje $\lambda_{m+n}(B\times D)=\lambda_m(B)\lambda_n(D)$
dla wszystkich borelowskich $B$ i $D$ ze
skończoną miarą drugiego czynnika. Dla dowolnego
borelowskiego $D$ wyczerpujemy go zbiorami
$D\cap[-k,k]^n$ i stosujemy ciągłość miary z dołu.
Tak otrzymujemy formułę także przy nieskończonych
miarach czynników, z umową $0\cdot\infty=0$.

Dla dowolnych $A,C$ jak wyżej wartość miary
prostokąta nie zmienia się po zastąpieniu go
odpowiednim prostokątem borelowskim, z wyjątkiem
przypadków nieskończonych miar, gdzie stosujemy
wyczerpanie obydwu czynników zbiorami ograniczonymi
i ciągłość miar z dołu. Zatem
$\lambda_{m+n}(A\times C)=\lambda_m(A)\lambda_n(C)$
(z konwencją dla iloczynu z zerem).
Jednoznaczność miary produktowej daje zgodność
miar na całej $\mathcal L_m\otimes\mathcal L_n$.

\emph{Krok 3: uzupełnienie.}
$\mathcal B(\mathbb R^{m+n})$ jest podzbiorem
$\mathcal L_m\otimes\mathcal L_n$, ponieważ
borelowska $\sigma$-algebra produktu euklidesowego
jest generowana przez prostokąty otwarte.
Każdy zbiór z $\mathcal L_{m+n}$ znajduje się
między dwoma zbiorami borelowskimi $B\subseteq E\subseteq C$
ze $\lambda_{m+n}(C\setminus B)=0$
(dzięki regularności miary Lebesgue'a i jej
zupełności). Zatem $E$ należy do uzupełnienia
$\mathcal L_m\otimes\mathcal L_n$.
W drugą stronę uzupełnienie pozostaje podzbiorem
$\mathcal L_{m+n}$, ponieważ ta ostatnia
$\sigma$-algebra jest zupełna. To kończy dowód.
\end{proof}

\begin{przyklad}[Niezmierzalny w produkcie podzbiór przekątnej]
\label{ex:diagonal-not-product-12}
Niech $V\subseteq[0,1]$ będzie zbiorem niemierzalnym
w sensie Lebesgue'a, np. zbiorem Vitaliego, i niech
$E=\{(x,x):x\in V\}\subseteq\mathbb R^2$.
Zbiór $E$ jest podzbiorem przekątnej
$D=\{(x,x):x\in[0,1]\}$, a $\lambda_2(D)=0$.
Zupełność $\lambda_2$ daje $E\in\mathcal L_2$.
Gdyby natomiast $E\in\mathcal L_1\otimes\mathcal L_1$,
to przeciwobraz $\Delta^{-1}(E)$ przez
$\Delta(x)=(x,x)$ byłby Lebesgue'a-mierzalny.
Rzeczywiście oba odwzorowania współrzędne
$\operatorname{pr}_j\circ\Delta$ są identycznością,
a więc $\Delta$ jest mierzalna jako odwzorowanie
do produktu przestrzeni mierzalnych. Jednak
$\Delta^{-1}(E)=V$, sprzeczność. Stąd
$\mathcal L_1\otimes\mathcal L_1\subsetneq\mathcal L_2$.
Co ciekawe, każdy pionowy i poziomy przekrój $E$
jest pusty lub jednoelementowy, a więc mierzalny:
to dodatkowo dowodzi, że nie można odwracać
lematu~\ref{lem:sections-measurable-12}.
\end{przyklad}

\begin{uwaga}[Pełne zbiory mierzalne i przekroje]
Z lematu~\ref{lem:sections-measurable-12} nie wynika,
że \emph{wszystkie} przekroje dowolnego zbioru
$E\in\mathcal L_{m+n}$ są mierzalne. Na przykład
dla niemierzalnego $V\subseteq[0,1]$ zbiór
$V\times\{0\}$ jest $\lambda_2$-mierzalny jako
podzbiór prostej poziomej miary zero, ale jego
przekrój przy $y=0$ jest równy $V$.
Wersja Fubiniego dla uzupełnionej miary mówi o
mierzalności i całkowalności przekrojów
\emph{prawie wszędzie}; do tego potrzebne jest
osobne uzasadnienie po zastąpieniu funkcji
przez mierzalnego przedstawiciela produktowego.
\end{uwaga}


\begin{twierdzenie}[Tonelli i Fubini dla uzupełnionej miary produktowej]
\label{thm:completed-product-fubini-12}
Niech $(X,\mathcal F,\mu)$ i $(Y,\mathcal G,\nu)$
będą przestrzeniami z miarami $\sigma$-skończonymi
i \emph{zupełnymi}. Niech $(\widehat{\mathcal H},
\widehat\pi)$ będzie uzupełnieniem przestrzeni
$(X\times Y,\mathcal F\otimes\mathcal G,\mu\otimes\nu)$.
Jeśli $f:X\times Y\to[0,\infty]$ jest
$\widehat{\mathcal H}$-mierzalna, to dla $\mu$-prawie
każdego $x$ przekrój $f_x$ jest $\mathcal G$-mierzalny,
a dla $\nu$-prawie każdego $y$ przekrój $f^y$ jest
$\mathcal F$-mierzalny. Po nadaniu wartości zero
na wyjątkowych mierzalnych zbiorach parametrów
funkcje powstałe przez całkowanie przekrojów są
mierzalne i spełniają wzór Tonellego z całką
względem $\widehat\pi$.

Jeśli $f$ jest rzeczywista lub zespolona,
$\widehat{\mathcal H}$-mierzalna i
$\int_{X\times Y}|f|\,\mathrm d\widehat\pi<\infty$,
to przekroje są całkowalne prawie wszędzie,
a wzór Fubiniego również zachodzi po pominięciu
zbiorów wyjątkowych.
\end{twierdzenie}
\begin{proof}
\emph{Krok 1: przedstawiciel mierzalny względem
$\mathcal F\otimes\mathcal G$.}
Z definicji uzupełnienia miary każda funkcja
$\widehat{\mathcal H}$-mierzalna o wartościach
w $\overline{\mathbb R}$ jest równa
$\widehat\pi$-prawie wszędzie pewnej funkcji
$g$ mierzalnej względem $\mathcal F\otimes\mathcal G$.
Aby to zobaczyć, dla $f\ge0$ wybieramy proste
mierzalne przybliżenia $s_n\uparrow f$, a każdy
zbiór poziomicowy każdej funkcji $s_n$ zastępujemy
zbiorem z $\mathcal F\otimes\mathcal G$ różniącym
się od niego tylko na podzbiorze mierzalnego
$\pi$-zbioru zerowego. Przeliczalna suma takich
zbiorów zerowych nadal ma $\pi$-miarę zero;
poza nią zastąpione funkcje proste mają tę samą
granicę co $f$. W przypadku funkcji ze znakiem
i zespolonych postępujemy osobno z częściami
rzeczywistymi i urojonymi. Daje to $g$ i zbiór
$N\in\mathcal F\otimes\mathcal G$ z $\pi(N)=0$
taki, że $f=g$ poza $N$.

\emph{Krok 2: przekroje zbioru zerowego.}
Tonelli dla $\chi_N$ daje
$0=\pi(N)=\int_X\nu(N_x)\,\mathrm d\mu(x)$.
Wobec tego $\nu(N_x)=0$ dla $\mu$-prawie każdego
$x$. Analogicznie $\mu(N^y)=0$ dla $\nu$-prawie
każdego $y$. Dla takich $x$ przekroje $f_x$ i $g_x$
różnią się jedynie na podzbiorze $N_x$.
Ponieważ $\nu$ jest zupełna, $f_x$ jest
$\mathcal G$-mierzalna i ma tę samą całkę co $g_x$;
podobnie z przekrojami poziomymi. Poza zbiorami
wyjątkowymi można zatem całkować $g$ zamiast $f$.

\emph{Krok 3: równość całek.}
Z równości $f=g$ $\widehat\pi$-prawie wszędzie
wynika równość ich całek względem $\widehat\pi$
w przypadku funkcji nieujemnej lub całkowalnej.
Na $g$ stosujemy odpowiednio twierdzenie
\ref{thm:tonelli-12} lub \ref{thm:fubini-12}.
Z kroku 2 całki przekrojowe $f$ i $g$ zgadzają
się prawie wszędzie, więc ich całki iterowane
są równe. To daje obie tezy.
\end{proof}

\section{Zbiory zerowe w produkcie i wyjątki przekrojowe}

\begin{wniosek}[Prawie wszystkie przekroje zbioru zerowego]
\label{cor:null-sections-12}
Niech $\mu,\nu$ będą miarami $\sigma$-skończonymi,
a $N\in\mathcal F\otimes\mathcal G$. Wówczas
$(\mu\otimes\nu)(N)=0$ wtedy i tylko wtedy, gdy
$\nu(N_x)=0$ dla $\mu$-prawie każdego $x$, a równoważnie
$\mu(N^y)=0$ dla $\nu$-prawie każdego $y$.
Nie wolno w tym stwierdzeniu zastępować
\emph{prawie każdego} przez \emph{każdego}.
\end{wniosek}
\begin{proof}
Na mocy twierdzenia~\ref{thm:section-identity-12}
obie funkcje przekrojowe są mierzalne i mają
nieujemne wartości. Z konstrukcji miary produktowej
\eqref{eq:measure-section-12} ich całki są równe
$(\mu\otimes\nu)(N)$. Całka funkcji nieujemnej
jest równa zeru dokładnie wtedy, gdy ta funkcja
znika prawie wszędzie. To daje obie równoważności.
Dla przykładu, gdy $X=Y=[0,1]$ z miarami Lebesgue'a,
zbiór $N=\{0\}\times[0,1]$ ma miarę produktową zero,
lecz $\lambda_1(N_0)=1$.
\end{proof}

\section{Więcej niż dwa czynniki i zastosowanie do prawdopodobieństwa}

\begin{wniosek}[Łączność skończonych produktów]
\label{cor:associative-product-12}
Niech $(X_j,\mathcal F_j,\mu_j)$, $j=1,2,3$,
będą przestrzeniami z miarami $\sigma$-skończonymi.
Przy naturalnej identyfikacji
$(X_1\times X_2)\times X_3=X_1\times(X_2\times X_3)$
zachodzi
$(\mu_1\otimes\mu_2)\otimes\mu_3
=\mu_1\otimes(\mu_2\otimes\mu_3)$.
W szczególności dla każdej nieujemnej funkcji
mierzalnej na iloczynie trzech przestrzeni
kolejność całkowania jest dowolna, a dla funkcji
bezwzględnie całkowalnej można stosować Fubiniego
w dowolnej kolejności.
\end{wniosek}
\begin{proof}
Oba konstrukty są miarami na tej samej
$\sigma$-algebrze produktowej, ponieważ
$\sigma$-algebra generowana przez trójwymiarowe
prostokąty nie zależy od umieszczenia nawiasów.
Dla każdego prostokąta $A_1\times A_2\times A_3$
obie miary przyjmują wartość
$\mu_1(A_1)\mu_2(A_2)\mu_3(A_3)$
(z konwencją dla zer i nieskończoności).
Miary są $\sigma$-skończone jako produkty miar
$\sigma$-skończonych. Jednoznaczność rozszerzenia
na generującym układzie prostokątów zapewnia ich
równość. Twierdzenia o całkach wynikają przez
kolejne zastosowanie Tonellego i Fubiniego.
\end{proof}

\begin{przyklad}[Niezależność dwóch zmiennych losowych]
Niech $P$ i $Q$ będą miarami probabilistycznymi
na $(X,\mathcal F)$ i $(Y,\mathcal G)$.
Na produkcie $(X\times Y,\mathcal F\otimes\mathcal G)$
rozważmy $P\otimes Q$ i odwzorowania współrzędne
$U:X\times Y\to X$, $U(x,y)=x$, oraz
$V:X\times Y\to Y$, $V(x,y)=y$.
Oba są mierzalne względem odpowiednich $\sigma$-algebr.
Dla $A\in\mathcal F$ i $B\in\mathcal G$
mamy
$(P\otimes Q)(\{U\in A,V\in B\})=P(A)Q(B)$.
To jest definicja niezależności zmiennych losowych
$U$ i $V$. Jeżeli $\varphi:X\to[0,\infty]$ oraz
$\psi:Y\to[0,\infty]$ są mierzalne, Tonelli daje
\[
\int_{X\times Y}\varphi(x)\psi(y)
\,\mathrm d(P\otimes Q)(x,y)
=\left(\int_X\varphi(x)\,\mathrm dP(x)\right)
\left(\int_Y\psi(y)\,\mathrm dQ(y)\right).
\]
Przy funkcjach zespolonych tożsamość obowiązuje,
jeżeli są całkowalne względem odpowiednich miar.
\end{przyklad}

\section{Podsumowanie i dalszy kierunek}
Najpierw udowodniliśmy mierzalność przekrojów,
a następnie \emph{bez użycia Fubiniego} wykazaliśmy
równość całek przekrojowych w prostokątach o
skończonej mierze. Twierdzenie o klasie monotonicznej
oraz wyczerpanie przez prostokąty $\sigma$-skończone
umożliwiły konstrukcję jednoznacznej miary
$\mu\otimes\nu$. Dopiero z tej konstrukcji
wyprowadziliśmy Tonellego i Fubiniego.

Warto zapamiętać rozdział ról: Tonelli pozwala
przenosić kolejność dla funkcji \emph{nieujemnej},
także o nieskończonej całce. Fubini zapewnia
pełną równość całek funkcji zmieniających znak,
lecz wymaga kontroli $\int|f|$. Przypadki
miar nie-$\sigma$-skończonych i niezupełność
$\sigma$-algebry produktowej ostrzegają przed
nadmiernym uogólnianiem tych wyników.
